\documentclass[conference]{IEEEtran}
\usepackage[utf8]{inputenc}
\usepackage[T1]{fontenc}
\usepackage[brazil]{babel}
\usepackage{amsmath}
\usepackage{booktabs}
\usepackage{listings}
\usepackage{xcolor}
\usepackage{multicol}
\usepackage{hyperref}

\lstdefinestyle{pla}{
  basicstyle=\ttfamily\small,
  frame=single,
  breaklines=true,
  columns=fullflexible,
  backgroundcolor=\color{gray!10},
  aboveskip=4pt,
  belowskip=4pt
}

\title{Algoritmo de Quine-McCluskey em C++ | Parte prática apenas}

\begin{document}

\maketitle

\begin{abstract}
Este artigo documenta a implementação prática do algoritmo de Quine-McCluskey
em C++, um minimizador de funções booleanas baseado em leitura de arquivos no
formato PLA. São apresentados a estrutura do programa, o formato de entrada e
saída adotado, e os resultados obtidos em três casos de teste.
\end{abstract}

\section{Introdução}

A minimização de funções booleanas é um problema central no projeto de circuitos
digitais. O algoritmo de Quine-McCluskey \cite{mccluskey1956} resolve esse
problema de forma sistemática para qualquer número de variáveis. A implementação
desenvolvida recebe um arquivo PLA como entrada, executa a minimização e gera
um arquivo PLA de saída com a expressão booleana simplificada \cite{stroustrup2013}.

\section{Estrutura do Programa}

O programa é executado via linha de comando, recebendo o caminho do arquivo
de entrada como argumento:

\begin{lstlisting}[style=pla]
./src input/arquivo.pla
\end{lstlisting}

Os arquivos de entrada ficam na pasta \texttt{input/} e o resultado é salvo
automaticamente em \texttt{output/result.pla}. O fluxo interno segue três
etapas: (1)~leitura do PLA e extração dos mintermos com saída 1;
(2)~geração dos implicantes primos por combinações sucessivas;
(3)~seleção da cobertura mínima.

\section{Formato PLA}

A diretiva \texttt{.i} indica o número de entradas e \texttt{.o} o número de
saídas. Cada linha de dados contém os bits de entrada, um espaço e o valor
da saída. O símbolo \texttt{-} representa \textit{don't care} e \texttt{.e}
encerra o arquivo.

\section{Resultados}

\subsection{Caso 1: Função AND de 2 variáveis}

\noindent\textbf{Entrada} (\texttt{input/and.pla}):
\begin{lstlisting}[style=pla]
.i 2
.o 1
11 1
.e
\end{lstlisting}

\noindent\textbf{Saída} (\texttt{output/result.pla}):
\begin{lstlisting}[style=pla]
F = AB

PLA output:
.i 2
.o 1
11 1
.e
\end{lstlisting}

A função AND só vale 1 quando ambas as entradas são 1. Como há apenas um
mintermo, nenhuma simplificação é possível e o programa retorna $F = AB$.

\subsection{Caso 2: Função de 3 variáveis}

\noindent\textbf{Entrada} (\texttt{input/ex01\_funcao3var.pla}):
\begin{lstlisting}[style=pla]
.i 3
.o 1
0-0 1
-11 1
.e
\end{lstlisting}

\noindent\textbf{Saída} (\texttt{output/result.pla}):
\begin{lstlisting}[style=pla]
F = A'C' + BC

PLA output:
.i 3
.o 1
0-0 1
-11 1
.e
\end{lstlisting}

O implicante \texttt{0-0} representa $A'C'$, com $B$ irrelevante, e
\texttt{-11} representa $BC$, com $A$ irrelevante.

\subsection{Caso 3: Função de 4 variáveis}

\noindent\textbf{Entrada} (\texttt{input/ex00\_f\_da\_folha.pla}):
\begin{lstlisting}[style=pla]
.i 4
.o 1
0-00 1
-000 1
.e
\end{lstlisting}

\noindent\textbf{Saída} (\texttt{output/result.pla}):
\begin{lstlisting}[style=pla]
F = A'C'D' + B'C'D'

PLA output:
.i 4
.o 1
0-00 1
-000 1
.e
\end{lstlisting}

O implicante \texttt{0-00} representa $A'C'D'$ e \texttt{-000} representa
$B'C'D'$, cobrindo exatamente os mintermos 0, 4 e 8.

\section{Código Fonte}

O código fonte completo da implementação está disponível publicamente em:

\begin{center}
\url{https://github.com/LuccaDultra/Relatorios-SD/tree/main/PSD/Ciclo\%202/Lab\%206/qcluskey-algorithm}
\end{center}

\section{Conclusão}

A implementação mostrou-se correta nos três casos testados, produzindo
expressões mínimas coerentes com os mintermos fornecidos. O uso do formato
PLA tanto na entrada quanto na saída facilita a integração com outras
ferramentas de síntese lógica.

\begin{thebibliography}{1}

\bibitem{mccluskey1956}
E.~J. McCluskey, ``Minimization of Boolean Functions,''
\textit{Bell System Technical Journal}, vol.~35, no.~6,
pp.~1417--1444, 1956.

\bibitem{stroustrup2013}
B.~Stroustrup, \textit{The C++ Programming Language}, 4th~ed.
Upper Saddle River, NJ: Addison-Wesley, 2013.

\end{thebibliography}

\end{document}